\section{Asymptotic distance and linear work}
\label{sec:structured-distance}

We now instantiate Structured SPIN with a growing outer and fixed IMT maps.
The outer length grows logarithmically, but its encoding work is linear.
The proof must handle both messages concentrated in a few outer blocks
and messages occupying a constant fraction of the blocks.

\subsection{Scalable constituent family}
\label{sec:structured-scalable}

The scalable outer begins with an even base length $a$ and an injective
rate-one-half map $H_a:\F_2^{a/2}\longrightarrow\F_2^a$.
Its weight enumerator
\[
  W_a(z):=\sum_{v\in\F_2^{a/2}}z^{\wt(H_a(v))}
\]
must be known exactly or bounded by a certified envelope. For any multiple
$b$ of $a$, let $H_{a,b}$ be the direct sum of $b/a$ copies of $H_a$.
Setup samples two independent permutations $\tau_1,\tau_2\gets S_b$ and
constructs the rate-one-half BA-3 constituent
\begin{equation}
  \mathcal O_{a,b}
  :=
  \Acc_b\circ\tau_2\circ
  \Acc_b\circ\tau_1\circ H_{a,b}.
  \label{eq:ba3-outer}
\end{equation}
It samples $\mathcal O_{a,b}$ once and reuses that realization in every outer
position. This is the chosen-block form of the BA construction
\cite{KolesnikovEtAl2026BA,cryptoeprint:2026/1903}.

The base length $a$ is a construction parameter. A distance certificate must
bind a particular $H_a$, its spectrum, and an admissible schedule for $b$.
Changing $a$ or $H_a$ changes the outer enumerator and may change the certified
schedule. It does not change the route or the structured inner.

\paragraph{Current certified choice.}
For the theorem below, set $a:=24$ and let
$H_{24}:\F_2^{12}\longrightarrow\F_2^{24}$ encode the extended binary Golay code. Fix a positive integer $m$ and set
$L_m:=128m$. Let $b_m$ be the least positive multiple of $24$ satisfying
\begin{equation}
  b_m\ge \frac{39}{4}\log_2(L_mb_m),
  \qquad
  N_m:=L_mb_m,
  \qquad
  K_m:=N_m/2.
  \label{eq:structured-native-schedule}
\end{equation}
Then $b_m=\frac{39}{4}\log_2N_m+O(1)$, and consecutive direct lengths have relative
gap $o(1)$.

Use $\mathcal O_{24,b_m}$ from \eqref{eq:ba3-outer}. Setup reuses this sampled
map in all $L_m$ diagonal positions:
$G_{m,1}=\cdots=G_{m,L_m}:=\mathcal O_{24,b_m}$.
The expectation over $(\tau_1,\tau_2)$ remains outside the
$q$th enumerator power, as in
\eqref{eq:shared-constituent-enumerator}.

After sampling this shared constituent, setup follows the generic
construction. It samples one independent $\sigma_j\gets S_{b_m}$ for every
outer position and one independent $\pi_r\gets S_{L_m}$ for every region.
The inner fixes $(t,s)=(128,19)$ and the IMT maps $A,C$ in
Table~\ref{tab:imt-maps}. Setup samples one independent transvection per
inner step as above. All route permutations and transvections are independent
of $(\tau_1,\tau_2)$.

The selected maps have a small, explicit weight structure. For nonzero
states $q$, the only values of $\wt(Aq)$ are $48,56,64,72,80$.
The transfer analysis tracks these five shells rather than all $2^{19}$
states. The $128$ columns of $C$ are distinct nonzero weight-five vectors
spanning $\F_2^{19}$, so each input coordinate contributes to only five
state coordinates. This gives sparse feedback without treating its
syndromes as uniform: the proof uses the exact spectrum of $C^\top$ to
bound cancellations. The appendix supplies the map bases, shell
multiplicities, and transfer entries; Section~\ref{sec:implementation}
describes the circuits used to apply the maps.

\begin{theorem}[Scalable Structured SPIN]
\label{thm:structured-spin-scalable}
For the setup above,
\begin{equation}
  \Prob\left[
    d_{\min}(E_m)\le\lfloor0.11N_m\rfloor
  \right]
  =o(1)
  \qquad(m\to\infty).
  \label{eq:structured-spin-distance}
\end{equation}
The probability includes the two BA permutations, all route permutations,
and all IMT transvections. Every realized code has rate $1/2$. Its ordinary
and transposed encoders use $O(N_m)$ bit operations.
\end{theorem}

The argument has three parts. First, control the realized spectrum of the
one sampled outer. Next, bound the route--inner failure probability for
each number of active blocks. Finally, sum those bounds before applying
the first-moment theorem. Appendix~\ref{app:structured-scalable-certificate}
gives the spectrum calculations, transfer entries, and numerical witnesses
used below.
The block constant $39/4$ is the current certified choice; improving it
would change the schedule, not the construction.

\subsection{Selecting the shared outer}

Suppress the family index $m$, and write $\mathcal O_b:=\mathcal O_{24,b}$.
Let $A_{\mathcal O_b}(w)$ count its nonzero words of weight $w$, and let
$\overline A_b(w):=\E_{\tau_1,\tau_2}A_{\mathcal O_b}(w)$.
The exact formula in Lemma~\ref{lem:structured-exact-ba-spectrum} composes
the Golay enumerator with two accumulator weight transitions.
It supplies both the tail bound and the spectrum envelope.
Let $\mathcal W:=[0.104,0.896]$. Lemma~\ref{lem:structured-ba-tails}
proves
\begin{equation}
  \sum_{w/b\notin\mathcal W}\overline A_b(w)=o(1).
  \label{eq:structured-ba-tail}
\end{equation}

\begin{lemma}[One-sample selection]
\label{lem:structured-one-sample-selection}
There is an event $\mathcal G_b$, determined by the two BA permutations, such
that
\begin{equation}
  \Prob[\mathcal G_b]=1-o(1).
  \label{eq:ba-good-event}
\end{equation}
On $\mathcal G_b$, every nonzero constituent word has relative weight in
$\mathcal W$, and
\begin{equation}
  A_{\mathcal O_b}(w)
  \le b^2\overline A_b(w)
  \qquad(0\le w\le b).
  \label{eq:ba-selected-spectrum}
\end{equation}
\end{lemma}

\begin{proof}
The expected number of nonzero tail words equals the left side of
\eqref{eq:structured-ba-tail}. Markov's inequality shows that no such word
exists with probability $1-o(1)$. For each remaining weight $w$, another
application of Markov's inequality gives
\[
  \Prob[A_{\mathcal O_b}(w)>b^2\overline A_b(w)]\le b^{-2}.
\]
A union bound over at most $b+1$ weights costs $O(1/b)$. Intersecting the two
events proves the lemma.
\end{proof}

The event in Lemma~\ref{lem:structured-one-sample-selection} is paid once.
After setup samples $\mathcal O_b$, the construction reuses that realization
at every outer position.

Fix any realization satisfying $\mathcal G_b$. All subsequent bounds are
uniform over this choice; the remaining randomness consists of the
independent route permutations and IMT transvections.
The resulting bounds therefore also hold conditional on $\mathcal G_b$.
The spectrum calculation also gives a concave function
$\widehat a_{\rm BA}$ on $\mathcal W$ such that
\[
 A_{\mathcal O_b}(w)
 \le \exp\{b\widehat a_{\rm BA}(w/b)+O(\ln b)\}.
\]
Appendix~\ref{app:structured-scalable-certificate} derives this majorant
from the exact expected spectrum, including the uniform remainder.
Selection lets us multiply bounds on the realized spectrum across rows.
It does not replace the expectation of a power by a power of the mean.

For $0\le d\le N$ and $1\le Q\le L$, let $Z_{d,Q}$ count nonzero messages
with exactly $Q$ active outer blocks and output weight at most $d$.
Our goal is to show that the sum of their conditional expectations is $o(1)$.

\subsection{From row weights to an iid comparison}

View an outer word as an $L$-tuple of rows in $\F_2^b$. For the framework in
Section~\ref{sec:framework}, take its class to be the ordered row-weight
profile $c(u):=(\wt(u_1),\ldots,\wt(u_L))$.
This finite class records the information preserved by the independent block
shuffles. The analysis below constructs a uniform route--inner failure
envelope for every such profile. It then aggregates profiles by occupation and
mean row weight where those coarser parameters suffice.

Fix $Q$ active outer positions. Let their realized output weights be
$w_1,\ldots,w_Q$, and set
\[
  x_j:=\frac{w_j}{b},
  \qquad
  \alpha:=\frac QL,
  \qquad
  x:=\frac1Q\sum_{j=1}^Qx_j,
  \qquad
  q:=\alpha x.
\]
The next lemma justifies the reduction from the complete weight tuple to
$(\alpha,x)$ when $\alpha$ has a positive lower bound.
We use binary relative entropy with natural logarithms,
\[
 D_{\rm KL}(u\Vert v)
 :=u\ln(u/v)+(1-u)\ln((1-u)/(1-v)),
\]
with the usual continuous endpoint conventions.

\begin{lemma}[Route domination]
\label{lem:structured-route-domination}
Let $F$ be any nonnegative function of the routed $L$-by-$b$ binary array.
Conditioned on the row weights $w_1,\ldots,w_Q$, the structured route
satisfies
\begin{equation}
  \E_{\rm route}[F]
  \le
  (b+1)^Q(L+1)^b
  \E_{\operatorname{Ber}(q)^{Lb}}[F].
  \label{eq:structured-route-domination}
\end{equation}
If $\alpha$ is bounded below by a positive constant, the multiplicative
factor in \eqref{eq:structured-route-domination} is $\exp(o(N))$.
\end{lemma}

\begin{proof}
For row $j$, begin with $b$ independent Bernoulli-$x_j$ variables and
condition their sum to equal $w_j$. The value $w_j$ is a mode of the
binomial distribution. Hence the inverse conditioning probability is at
most $b+1$. Applying this reduction to all active rows costs at most
$(b+1)^Q$.

Before the region shuffle, one fixed region contains independent Bernoulli
variables whose average success probability is $q$. Let $S$ be their sum.
For every $0\le k\le L$, the Chernoff bound for a Poisson-binomial sum gives
\[
  \Prob[S=k]
  \le
  \exp\{-L D_{\rm KL}(k/L\Vert q)\}.
\]
The method-of-types lower bound gives
\[
  \Prob[\operatorname{Bin}(L,q)=k]
  \ge
  \frac1{L+1}
  \exp\{-L D_{\rm KL}(k/L\Vert q)\}.
\]
After a uniform region shuffle, both laws are uniform on every weight-$k$
slice. The routed region law is therefore pointwise at most $L+1$ times the
iid Bernoulli-$q$ law. Multiplying over the $b$ independent region shuffles
proves \eqref{eq:structured-route-domination}.

Finally,
\[
  \ln((b+1)^Q(L+1)^b)
  =Q\ln(b+1)+b\ln(L+1)=o(Lb)
\]
when $Q/L$ is bounded below by a positive constant and $b,L\to\infty$.
\end{proof}

Concavity of $\widehat a_{\rm BA}$ now reduces the selected outer spectrum
to the mean row weight:
\begin{equation}
  \frac1Q\sum_{j=1}^Q\widehat a_{\rm BA}(x_j)
  \le
  \widehat a_{\rm BA}(x).
  \label{eq:structured-ba-jensen}
\end{equation}
The region shuffles are essential for Lemma~\ref{lem:structured-route-domination}.
The deterministic transpose alone spreads each row, but it does not erase
the row's coordinate pattern.

\subsection{IMT moments at positive occupation}

The iid comparison is useful because the inner has a finite-state
description. For iid Bernoulli-$\beta$ input and $0<z<1$, weight each
transition by $z$ to the emitted Hamming weight. Track mass at zero,
arbitrary nonzero mass, and uniform mass on each of the five nonzero
image-weight shells of $A$. This gives seven nonnegative coordinates.
The lazy branch can cancel through $CX=Q_i$; the refresh branch spreads
a live state uniformly. Both effects enter the transfer bound.

Appendix~\ref{app:imt-finite-transfers} constructs nonnegative matrices
$T(\beta,z)$ that dominate these weighted measures step by step.
The occupation envelope $T_{\rm occ}(\beta,z)$ averages the transfer bounds
for uniform fixed-weight inputs with weights from $\operatorname{Bin}(128,\beta)$.
Starting with the zero-state unit row vector $e_Z$, induction gives
\[
 \E[z^{\wt(Y)}]\le e_Z T(\beta,z)^{N/128}\1,
\]
where $\1$ is the all-ones column vector.
A positive column vector $\boldsymbol v$ satisfying
$T\boldsymbol v\le\lambda\boldsymbol v$ therefore gives a bound
$C_{\boldsymbol v}\lambda^{N/128}$, with
$C_{\boldsymbol v}:=v_Z/\min_i v_i$.
This is a componentwise inequality, not an assumption that numerical
eigenvalues are exact. The appendix also gives a scalar envelope.

Fix $\alpha\ge10^{-4}$ and mean row weight $x\in\mathcal W$.
Markov's inequality contributes $z^{-0.11N}$.
The outer spectrum and concavity contribute
$\exp\{N\alpha\widehat a_{\rm BA}(x)+o(N)\}$.
The factors for choosing active positions, row weights, and routing are
also $\exp(o(N))$, uniformly in this range.

We may choose a reference density $\beta=py$ rather than $\alpha x$.
Every routed array in the profile has weight $N\alpha x$.
On that slice, changing the reference law costs
$\exp\{ND_{\rm KL}(\alpha x\Vert py)\}$.
The chain rule bounds this divergence by
$D_{\rm KL}(\alpha\Vert p)+\alpha D_{\rm KL}(x\Vert y)$.
The resulting exponent per output bit is
\[
 \alpha\widehat a_{\rm BA}(x)
 +D_{\rm KL}(\alpha\Vert p)+\alpha D_{\rm KL}(x\Vert y)
 +\frac{\ln\lambda}{128}-0.11\ln z.
\]
The certificate fixes witnesses on finitely many boxes covering
$[10^{-4},1]\times\mathcal W$. On each affine segment of the outer majorant,
this expression is convex in $(\alpha,\alpha x)$; vertex bounds suffice.
Outward arithmetic gives an upper bound below $-4.10335\cdot10^{-7}$.
The finite set of witnesses has bounded prefactors, so the total
positive-occupation contribution is at most $L\exp(-\eta N+o(N))$
for some $\eta>0$.

\subsection{Sparse occupation without a large conditioning loss}

The dense comparison loses a factor $(L+1)^b$.
Although subexponential in $N$, this factor is too large when the useful
decay is only exponential in $Qb$. Sparse inputs need a different reference
measure.

Write $h(x):=-x\ln x-(1-x)\ln(1-x)$.
The outer majorant satisfies
$\widehat a_{\rm BA}(x)-h(x)+(\ln2)/2<\ell$ on $\mathcal W$,
with $\ell:=0.01281$.
After summing messages and shuffling each active row, comparison with
a fair $b$-bit row thus costs at most
$\exp\{b((\ln2)/2+\ell+\epsilon_b)\}$ per row,
where $\epsilon_b=O(\ln b/b)$.
This comparison concerns the message counting measure, not the law of a
single fixed word.

Under the fair-row reference, a region contains $Q$ uniformly placed
marks carrying independent fair bits. Generate marks independently with
probability $p=(8/5)\alpha$, then condition their count to equal $Q$.
Before conditioning, the bits have density $\beta=(4/5)\alpha$.
The inverse conditioning probability per region is at most
$8\sqrt Q\,\exp\{LD_{\rm KL}(\alpha\Vert p)\}$.
For this occupation envelope and $z=1-(8/5)\alpha$, an exact polynomial
check gives
\[
 T_{\rm occ}(\beta,z)\boldsymbol v(\alpha)
 \le(1-96\alpha)\boldsymbol v(\alpha)
 \qquad(0<\alpha\le10^{-4}),
\]
with moment prefactor at most $2048$.
Multiplying the row-counting, conditioning, and moment bounds yields
\[
 \begin{aligned}
 \E[Z_{\lfloor0.11N\rfloor,Q}\mid\mathcal G_b]
 \le{}&2048\binom LQ
 e^{Qb((\ln2)/2+\ell+\epsilon_b)}(8\sqrt Q)^b\\
 &{}\cdot e^{ND_{\rm KL}(\alpha\Vert p)}
 (1-96\alpha)^{N/128}z^{-0.11N}.
 \end{aligned}
\]
The schedule pays for choosing active positions, since
$\ln L/b\le4(\ln2)/39$. The leading coefficient per $Qb$ is below
$-0.01340384$. When $Q\ge4096$, the remaining conditioning cost per $Qb$
is below $0.002$.
The other remainders vanish uniformly as $b$ grows.
Consequently, for all sufficiently large lengths,
\[
 \E[Z_{\lfloor0.11N\rfloor,Q}\mid\mathcal G_b]
 \le e^{-0.006Qb},
 \qquad4096\le Q\le10^{-4}L.
\]
Appendix~\ref{app:imt-sparse} gives the polynomial witness and every remainder.
The explicit cutoff makes the bound uniform, even when $Q$ grows slowly.

\subsection{Fixed occupation: mixing between isolated impulses}

It remains to handle the finite set $Q<4096$.
Each region then contains at most $Q$ active bits.
As $L$ grows, those bits typically lie in distinct steps separated by
long empty intervals. A live state mixes during an empty interval:
after $g$ steps its law is
$2^{-g}$ times its entering law plus $(1-2^{-g})$ times the uniform
nonzero law. Mixing does not make the state zero.

Put $M:=2^{19}-1$ and $p_0:=2^{18}/M$.
Every coordinate of $A$ is nonzero, so a uniform live state emits
mean weight $128p_0$ per empty step.
The geometric mixing law bounds the variance over $g$ empty steps by
$3t^2g$. With the tilt $z_L=e^{-\theta/L}$, the weighted empty-interval
kernel therefore converges to
$D_\gamma(u):=\operatorname{diag}(1,e^{-\gamma u})$ on zero/live states,
where $\gamma=p_0\theta$ and $u$ is the interval's fraction of a region.

An isolated input bit starts a zero state because the columns of $C$
are nonzero. From a uniform live state, it cancels with probability $1/M$.
Its limiting transition is thus
\[
 P_0=\begin{pmatrix}0&1\\1/M&1-1/M\end{pmatrix}.
\]
For $a$ uniformly placed impulses, the limiting weighted region kernel is
\[
 K_a(\theta)=a!\int_{\substack{u_i\ge0\\\sum_{i=0}^a u_i=1}}
 D_\gamma(u_0)P_0D_\gamma(u_1)\cdots P_0D_\gamma(u_a)
 \,d\boldsymbol u.
\]
The integral uses $a$ independent coordinates.
Collisions and short gaps have probability $O_Q(\log L/L)$.
The empty-interval estimate makes the remaining error
$O_Q(L^{-1/2}+\log L/L)$.
Appendix~\ref{app:imt-fixed} proves a componentwise multiplicative comparison,
including the entries that vanish exactly. Across $b=O(\log L)$ regions,
its cost is $\exp(o(1))$.

For $Q=1,2$, combine these two-state kernels with the fair-row comparison.
The resulting exponents per $Qb$, including active-position choices,
are below $-0.10918$ and $-0.10914$, respectively.
For $Q\ge3$, retain each row weight: assign a positive variable $u_j$
to row $j$ and form the matrix polynomial
\[
 T_Q(\boldsymbol u,\theta)
 :=\sum_{\boldsymbol e\in\{0,1\}^Q}
 \left(\prod_j u_j^{e_j}\right)K^+_{\sum_j e_j}(\theta).
\]
Here $K_a^+$ replaces the live-to-live entry of $P_0$ by $1$, giving an
entrywise upper bound. Nonnegative coefficient extraction bounds the
fixed-row-weight moment by
\[
 \left(\prod_j\binom b{w_j}^{-1}u_j^{-w_j}\right)
 e_0^\top T_Q(\boldsymbol u,\theta)^b\1_2\,\exp(o_Q(b)).
\]
A weighted matrix norm separates this expression into one factor per row;
the outer majorant then bounds each factor.
The appendix proves that separation and verifies its endpoint inequalities.
For every fixed $3\le Q<4096$, it gives
\[
 \E[Z_{\lfloor0.11N\rfloor,Q}\mid\mathcal G_b]
 \le\exp\{-0.0008679Qb+o_Q(b)\}.
\]
The remainder may depend on $Q$: only finitely many such bounds are needed.

\subsection{Completing the argument}

\begin{proof}[Theorem~\ref{thm:structured-spin-scalable}]
Condition on any outer satisfying $\mathcal G_b$. The fixed-occupation
arguments give a vanishing sum over the finite set $1\le Q<4096$.
Equation~\eqref{eq:imt-sparse-final} gives a uniform geometric sum for
$4096\le Q\le10^{-4}L$. The positive-occupation certificate bounds the
remaining sum by $L\exp(-\eta N+o(N))$ for some $\eta>0$.
Thus $\sum_{Q=1}^L\E[Z_{\lfloor0.11N\rfloor,Q}\mid\mathcal G_b]=o(1)$.
The explicit cutoff makes this a summable union, not merely a statement
about each occupation sequence. Markov's inequality and
$\Prob[\mathcal G_b^c]=o(1)$ prove the distance claim.

Rate follows from injectivity. Every outer row costs $O(b)$ in either circuit
direction: it uses fixed Golay circuits, two permutations, and two accumulators.
The route costs $O(N)$, and the fixed-size binary IMT steps cost $O(1)$ each.
The displayed adjoint gives the same bound for transposed encoding.
\end{proof}

For any sufficiently large requested output length $n$, let $m(n)$ be the
largest index satisfying $N_{m(n)}\le n$. Define
\begin{equation}
  \widetilde E_n(x)
  :=
  E_{m(n)}(x)\Vert 0^{n-N_{m(n)}}.
  \label{eq:structured-requested-length-wrapper}
\end{equation}
This zero extension preserves injectivity and minimum distance. Consecutive
direct lengths have relative gap $o(1)$, so
$N_{m(n)}/n=1-o(1)$. The wrapped family therefore has rate
$1/2-o(1)$, relative distance greater than $0.11-o(1)$ with high probability,
and $O(n)$ ordinary and transposed work. The native-length theorem remains
the primary statement.
